% LaTeX preamble for the ATS PDF (create_pdf_from_md.py).
%
% Two goals: parse cleanly in ATS software, and let a human track the structure
% at a glance.
%
% Design rules used here:
%   1. Hierarchy comes from SIZE and SPACE, never from fading text out. A
%      heading is never lighter than the body text under it.
%   2. Colour is an accent only — the name, section titles and links. All body
%      copy and every job/project title stays near-black.
%   3. Grey is reserved for genuine metadata (dates), which is not a heading.
%   4. Single column, real text, no tables or images for layout, so ATS parsers
%      read it correctly. Colour itself is ignored by parsers and is safe.
%
% Pandoc maps the markdown as:
%   # Name       -> \section        ## Section -> \subsection
%   ### Company  -> \subsubsection  #### Project -> \paragraph
% In the article class the last two both default to bold \normalsize, which is
% why companies and projects were indistinguishable. Redefined below with base
% LaTeX \@startsection, so no extra TeX packages are needed (keeps CI light).

\usepackage{xcolor}

% Teal matches the site theme colour on saman.is-a.dev.
\definecolor{cvaccent}{HTML}{107878}
\definecolor{cvdark}{HTML}{15202B}
\definecolor{cvmeta}{HTML}{5A6675}

\linespread{1.05}
\setlength{\parskip}{0pt}
\color{cvdark}

% Left-align body text. LaTeX justifies by default, which stretches the spaces
% between words to force a flush right edge; on a dense CV that produces visibly
% uneven gaps line to line. \raggedright keeps word spacing constant.
% \setlength{\parindent}{0pt} because ragged text must not also indent.
\raggedright
\setlength{\parindent}{0pt}

% Accent rule drawn above each major section.
\newcommand{\cvrule}{%
  \par\vspace{0.2ex}%
  {\color{cvaccent}\hrule height 1pt}%
  \vspace{1.1ex}%
}

\makeatletter

% # Name — left aligned, large, accent colour.
\renewcommand\section{%
  \@startsection{section}{1}{0pt}%
    {0pt}%
    {0.5ex}%
    {\normalfont\huge\bfseries\color{cvaccent}}%
}

% ## Professional Summary / Work Experience — accent, rule above, one size up
% from a company heading so the two accent levels never collide.
\renewcommand\subsection{%
  \@startsection{subsection}{2}{0pt}%
    {3.0ex \@plus .4ex \@minus .2ex}%
    {1.1ex \@plus .2ex}%
    {\normalfont\Large\bfseries\color{cvaccent}\cvrule}%
}

% ### Company - dates — biggest item in the body, in the accent colour so an
% employer is never mistaken for one of the projects nested under it. Measured
% from the built PDF: company and project were both 12pt/10pt in the SAME
% near-black, and 2pt alone does not read as a level change at a glance.
\renewcommand\subsubsection{%
  \@startsection{subsubsection}{3}{0pt}%
    {2.8ex \@plus .4ex \@minus .2ex}%
    {0.8ex \@plus .2ex}%
    {\normalfont\large\bfseries\color{cvaccent}}%
}

% #### Project name — near-black, one step smaller than the company above it.
% Subordinate through colour AND size: the company is accent-coloured, the
% project is not, so the two levels separate at a glance.
\renewcommand\paragraph{%
  \@startsection{paragraph}{4}{0pt}%
    {1.6ex \@plus .2ex \@minus .1ex}%
    {0.35ex \@plus .1ex}%
    {\normalfont\normalsize\bfseries\color{cvdark}}%
}

\makeatother

% Links in the accent colour (applied late so it beats the pandoc template).
\AtBeginDocument{\hypersetup{urlcolor=cvaccent,linkcolor=cvaccent}}
